\section{第18章 工学基礎}
\hrule
\subsubsection*{この資料の主張}
ソフトウェア工学は、工学としての共通基盤を持つ。真の問題を見極め、制約の中で複数の解を比較し、証拠に基づいて判断し、結果を測定して改善することが重要である。第18章は、抽象化、実験、統計、測定、標準、根本原因分析を通じて、ソフトウェア技術者が工学的に判断するための土台を整理する。
\begin{notebox}{この資料の読み方}専門用語は原則として日本語で示し、必要に応じて英語を括弧内に併記する。英語名は、元資料やツール上の名称を確認したいときの手がかりとして残す。初学者が前提知識なしに読めるように、各節では「何のための概念か」「実務では何を見ればよいか」を重視する。\end{notebox}
